\section{Meta-FBOHN v1/v2: Negative Result of Feature Scaling}
\label{sec:meta_fbohn_v1v2}

\subsection{Hypothesis}
The first versions of Meta-FBOHN attempted to improve FBOHN by learning weights:
\[
x_i' = w_i x_i.
\]
Variant v1 used five block weights $(E,M,A,S,L)$, and variant v2 used 30~feature weights.

\subsection{Result}
For all seeds and noise levels, practically identical values were obtained:
\[
\text{Meta-FBOHN-v1}=\text{Meta-FBOHN-v2}=\text{FBOHN}.
\]

\subsection{Explanation}
After computing features, the classifier applied
\[
\texttt{StandardScaler}+\texttt{LogisticRegression}.
\]
For positive diagonal scaling $D=\mathrm{diag}(w_i)$, standardization removes the scaling effect:
\[
\frac{w_ix_i-w_i\mu_i}{w_i\sigma_i}=\frac{x_i-\mu_i}{\sigma_i}.
\]

\subsection{Conclusion}
This is not a failure of Meta-FBOHN, but an important negative result:
\[
\boxed{\text{Mere weighting of FBOHN features does not create new information.}}
\]
Meta-FBOHN must create new features, not merely rescale existing ones.
